% ===== 《深入理解 AI Agent》 номын ElegantBook-ийн тусгай тохиргоо =====
% Үндсэн тодотгох өнгө ElegantBook-ийн cyan сэдвээс ирнэ: \structurecolor = RGB(31,186,190)

% ── Текст ба математикийн фонт: зохиогчийн өгүүлэлтэй нийцсэн Palatino хэв маяг ──
% Өгүүллүүдийн урьдчилсан хэвлэлийн (arxiv.sty) загвар нь үндсэн бичвэр,
% математикт mathpazo (Palatino), серифгүйд helvet[scaled=.90] ашигладаг.
% Энд XeLaTeX-д тохирох хувилбаруудыг ашиглана: бичвэрт TeX Gyre Pagella
% (Palatino), математикт TeX Gyre Pagella Math, серифгүйд TeX Gyre Heros
% (Helvetica-гийн хувилбар). Үүнийг class-д даатгалгүй энд заавал тохируул:
% pandoc-ийн template class ажилласны дараа фонтыг Latin Modern болгож өөрчилдөг
% тул зөвхөн preamble түвшний тохиргоо үйлчилнэ. Class-ийн адил файлын нэрээр
% ачаална: macOS fontconfig texmf-dist-ийг индексжүүлдэггүй тул бүлгийн нэрээр олдохгүй.
\usepackage{unicode-math}
\setmainfont{texgyrepagella}[
  UprightFont = *-regular,
  BoldFont = *-bold,
  ItalicFont = *-italic,
  BoldItalicFont = *-bolditalic,
  Extension = .otf,
  Scale = 1.0]
\setsansfont{texgyreheros}[
  UprightFont = *-regular,
  BoldFont = *-bold,
  ItalicFont = *-italic,
  BoldItalicFont = *-bolditalic,
  Extension = .otf,
  Scale = 0.9]
\setmathfont{texgyrepagella-math}[Extension = .otf]

% ── Unicode тэмдэгтийн нөөц фонт (текст/моно фонтууд эдгээрийг агуулахгүй) ──
\usepackage{newunicodechar}
% macOS-д Arial Unicode MS / Heiti SC бий; Linux-д Apple-ийн эдгээр фонт байхгүй
% тул бүтээлтийг хэвийн явуулахын тулд олон тэмдэг дэмждэг чөлөөт фонтоор орлуулна.
\IfFontExistsTF{Arial Unicode MS}{\newfontfamily\fallbackfont{Arial Unicode MS}}{\newfontfamily\fallbackfont{DejaVu Sans}}
\IfFontExistsTF{Heiti SC}{\newunicodechar{★}{{\fontspec{Heiti SC}★}}}{\newunicodechar{★}{{\fontspec{Noto Sans CJK SC}★}}}
\newunicodechar{✓}{{\fallbackfont ✓}}
\newunicodechar{✗}{{\fallbackfont ✗}}
\newunicodechar{₀}{{\fallbackfont ₀}}
\newunicodechar{₁}{{\fallbackfont ₁}}
\newunicodechar{₂}{{\fallbackfont ₂}}
\newunicodechar{₃}{{\fallbackfont ₃}}
\newunicodechar{₄}{{\fallbackfont ₄}}
\newunicodechar{₅}{{\fallbackfont ₅}}
\newunicodechar{₆}{{\fallbackfont ₆}}
\newunicodechar{₇}{{\fallbackfont ₇}}
\newunicodechar{₈}{{\fallbackfont ₈}}
\newunicodechar{₉}{{\fallbackfont ₉}}
\newunicodechar{ⓐ}{{\fallbackfont ⓐ}}
\newunicodechar{ⓑ}{{\fallbackfont ⓑ}}
\newunicodechar{ⓒ}{{\fallbackfont ⓒ}}
\newunicodechar{ⓓ}{{\fallbackfont ⓓ}}
\newunicodechar{ⓔ}{{\fallbackfont ⓔ}}
% Латин бичвэрийн фонтод байхгүй математик / тусгай тэмдэгтүүд
\newunicodechar{≈}{{\fallbackfont ≈}}
\newunicodechar{♠}{{\fallbackfont ♠}}
\newunicodechar{♣}{{\fallbackfont ♣}}
\newunicodechar{♥}{{\fallbackfont ♥}}
\newunicodechar{♦}{{\fallbackfont ♦}}
% Бичвэр доторх грек үсгүүд (энд Latin Modern тэдгээрийг дэмжихгүй)
\newunicodechar{α}{{\fallbackfont α}}
\newunicodechar{β}{{\fallbackfont β}}
\newunicodechar{γ}{{\fallbackfont γ}}
\newunicodechar{δ}{{\fallbackfont δ}}
\newunicodechar{ε}{{\fallbackfont ε}}
\newunicodechar{η}{{\fallbackfont η}}
\newunicodechar{θ}{{\fallbackfont θ}}
\newunicodechar{λ}{{\fallbackfont λ}}
\newunicodechar{μ}{{\fallbackfont μ}}
\newunicodechar{π}{{\fallbackfont π}}
\newunicodechar{ρ}{{\fallbackfont ρ}}
\newunicodechar{σ}{{\fallbackfont σ}}
\newunicodechar{τ}{{\fallbackfont τ}}
\newunicodechar{φ}{{\fallbackfont φ}}
\newunicodechar{ω}{{\fallbackfont ω}}

% ── Үлдсэн CJK тэмдэгтүүд ────────────────────────────────────────
% Энэ нь англи хэвлэл боловч хятад үсгийн цөөн жишээ текст дотор үлдсэн:
% хятад Token болгон хуваах, орчуулга, яриа таних зэрэг тохиолдлууд.
% ctex/xeCJK-ийн хятад бичвэрийн бүх давхаргыг ачаалалгүйгээр зөвхөн эдгээр
% код цэгийг CJK фонтоор дүрсэлнэ. macOS: Heiti SC;
% Linux: Noto Sans CJK SC.
\IfFontExistsTF{Heiti SC}{\newfontfamily\cjkfallback{Heiti SC}}{\newfontfamily\cjkfallback{Noto Sans CJK SC}}
\newunicodechar{医}{{\cjkfallback 医}}
\newunicodechar{图}{{\cjkfallback 图}}
\newunicodechar{型}{{\cjkfallback 型}}
\newunicodechar{模}{{\cjkfallback 模}}
\newunicodechar{生}{{\cjkfallback 生}}
\newunicodechar{解}{{\cjkfallback 解}}
\newunicodechar{陈}{{\cjkfallback 陈}}
\newunicodechar{两}{{\cjkfallback 两}}
\newunicodechar{么}{{\cjkfallback 么}}
\newunicodechar{了}{{\cjkfallback 了}}
\newunicodechar{京}{{\cjkfallback 京}}
\newunicodechar{代}{{\cjkfallback 代}}
\newunicodechar{体}{{\cjkfallback 体}}
\newunicodechar{北}{{\cjkfallback 北}}
\newunicodechar{右}{{\cjkfallback 右}}
\newunicodechar{大}{{\cjkfallback 大}}
\newunicodechar{天}{{\cjkfallback 天}}
\newunicodechar{左}{{\cjkfallback 左}}
\newunicodechar{怎}{{\cjkfallback 怎}}
\newunicodechar{思}{{\cjkfallback 思}}
\newunicodechar{推}{{\cjkfallback 推}}
\newunicodechar{智}{{\cjkfallback 智}}
\newunicodechar{样}{{\cjkfallback 样}}
\newunicodechar{概}{{\cjkfallback 概}}
\newunicodechar{气}{{\cjkfallback 气}}
\newunicodechar{点}{{\cjkfallback 点}}
\newunicodechar{理}{{\cjkfallback 理}}
\newunicodechar{的}{{\cjkfallback 的}}
\newunicodechar{考}{{\cjkfallback 考}}
\newunicodechar{能}{{\cjkfallback 能}}
\newunicodechar{零}{{\cjkfallback 零}}

% ── Код/тогтмол өргөнтэй фонт (код дахь хайрцаг зурах тэмдэгтүүдийг дэмжинэ) ──
% Menlo бол macOS-ийн фонт; Linux-д DejaVu Sans Mono-г нөөцөөр ашиглана.
\IfFontExistsTF{Menlo}{\setmonofont[Scale=0.86]{Menlo}}{\setmonofont[Scale=0.86]{DejaVu Sans Mono}}

% ── tcolorbox (ElegantBook аль хэдийн ачаалсан; зөвхөн Library нэмнэ) ──
\tcbuselibrary{skins,breakable}

% ── TikZ (cover.tex дэх зах хүртэл дүүргэсэн хавтасны зурагт) ─────
\usepackage{tikz}
\usetikzlibrary{fadings}
% Босоо бүдгэрэлт: дээд хэсэгт үл үзэгдэж, доод хэсэгт тодорно;
% ингэснээр хавтасны зураг хар хөх гарчгийн зурваст жигд уусна.
\tikzfading[name=titlefade, top color=transparent!100, bottom color=transparent!0]

% Сэдвийн тодотгох өнгө: ElegantBook-ийн тод cyan-г гүн хар хөхөөр солино
% (сонгодог, итгэл төрүүлэх техникийн номын өнгө аяс). Гарчиг, шугам, холбоос,
% кодын хүрээ, хавтсыг удирдана. Энэ мөрийг өөрчилж номын өнгийг бүхэлд нь шинэчил.
\definecolor{structurecolor}{RGB}{30,58,107}

% Код/зургийн хүрээний тодотгох өнгө (дээрх сэдвийн өнгийг дагана)
\colorlet{accent}{structurecolor}

% Дотоод харилцан эшлэлийн холбоос (图N-M / 第N章) — товшигдох, бүдэг ногоовтор,
% class-ийн нийтлэг linkcolor-оос тусдаа. crossref.lua ашиглана.
\newcommand{\crossreflink}[2]{\hyperref[#1]{\textcolor{structurecolor}{#2}}}
\colorlet{accentbg}{structurecolor!6}
\colorlet{accentrule}{structurecolor!35}

% ── Кодын блок дахь урт мөрийг таслах ───────────────────────────
% pandoc нь ```lang блокт `Highlighting`-ээр fancyvrb-ийн Verbatim,
% хэл заагаагүй блокт энгийн `verbatim` ашиглана; аль аль нь анхдагчаар
% мөр таслахгүй тул урт мөр хуудаснаас халин гарна. fvextra нь breaklines нэмнэ;
% хоёр орчныг мөр таслахаар дахин тодорхойлж, зайгүй хэсгийг Token дундаас
% тасална (URL, урт таних тэмдэг, CJK тайлбар). Үргэлжилсэн мөр бүрийг ↳ тэмдэглэнэ.
\usepackage{fvextra}
\fvset{breaklines=true,breakanywhere=true,%
  breaksymbolleft={\footnotesize\textcolor{accentrule}{$\hookrightarrow$}\,}}
\DefineVerbatimEnvironment{Highlighting}{Verbatim}{commandchars=\\\{\},breaklines,breakanywhere}

% ── Кодын блок: O’Reilly хэв маягийн нарийн хүрээтэй хайрцаг ───────
% БҮХ кодын блокт нэг хэв маяг хэрэглэж, өнгөөр ялгасан (```lang → Shaded)
% болон энгийн (``` → verbatim) блокийг адил харагдуулна. Энгийн verbatim
% хайрцагт fvextra Verbatim агуулна; \VerbatimEnvironment нь орчныг
% \end{Verbatim} биш \end{verbatim} (гадаад орчны нэр)-д төгсгөнө.
\newtcolorbox{codebox}{%
  breakable, enhanced,
  colback=black!3, colframe=accentrule,
  boxrule=0.5pt, leftrule=2.5pt,
  left=0.7em, right=0.6em, top=0.4em, bottom=0.4em,
  arc=0.6mm,
  before skip=0.7em, after skip=0.7em
}
\renewenvironment{Shaded}{\begin{codebox}}{\end{codebox}}
\renewenvironment{verbatim}%
  {\VerbatimEnvironment\begin{codebox}\begin{Verbatim}}%
  {\end{Verbatim}\end{codebox}}

% ── Туршилт ба асуултын тусгай хайрцаг (experiment_box.lua ашиглана) ──
\newtcolorbox{experimentbox}{
  breakable, enhanced,
  colback=accentbg, colframe=accent,
  boxrule=0.6pt, left=0.8em, right=0.8em, top=0.6em, bottom=0.6em,
  arc=1mm, before skip=0.9em, after skip=0.9em
}

\newtcolorbox{questionbox}{
  breakable, enhanced,
  colback=accentbg, colframe=accent,
  boxrule=0.8pt, left=0.9em, right=0.9em, top=0.9em, bottom=0.7em,
  arc=1.2mm, before skip=0.9em, after skip=0.9em,
  title={\sffamily\bfseries Гүнзгий бодох},
  fonttitle=\normalsize\color{white},
  coltitle=white,
  attach boxed title to top left={xshift=1em, yshift=-\tcboxedtitleheight/2},
  boxed title style={colback=accent, colframe=accent, arc=0.8mm, outer arc=0.8mm}
}

% ── Зургууд: эшлэл/асуултын хайрцагт багтахын тулд [H]-г албадна ──
\usepackage{float}
\let\origfigure\figure
\let\origendfigure\endfigure
\renewenvironment{figure}[1][htbp]{\origfigure[H]}{\origendfigure}
\setkeys{Gin}{width=\linewidth,height=0.38\textheight,keepaspectratio}

% ── Зургийн тайлбар: дугаарыг текстээр өгдөг тул автоматаар нэмэхгүй ──
\usepackage{caption}
% Серифгүй, тод, төвлөрсөн. Зургийн дугаар тайлбарын текстэд өөрт нь байна
% (crossref.lua), тиймээс автомат «Figure N:» шошгыг дарна.
\captionsetup{font={normalsize,sf,bf}, labelformat=empty, skip=4pt, justification=centering}

% ── Эшлэлийн блок: зүүн талын ногоовтор зурвас (туршилт эшлэлээр дүрслэгдэнэ) ──
\renewenvironment{quote}{%
  \begin{tcolorbox}[
    blanker, breakable,
    borderline west={2.5pt}{0pt}{accent},
    colback=accent!4,
    left=1em, right=0.6em, top=0.5em, bottom=0.5em,
    before skip=0.6em, after skip=0.6em
  ]%
}{%
  \end{tcolorbox}%
}

% ── Inline кодыг авсаархан байлгана ─────────────────────────────
\setlength{\parindent}{2em}

% ElegantBook нь бүлгийн эхний хуудсанд ашиглагдах `plain` хуудасны хэв маягийг хоосолдог.
% Энгийн хуудасны төвийн доод дугаартай адил эдгээр хуудсыг мөн дугаарлана.
\fancypagestyle{plain}{%
  \fancyhf{}%
  \fancyfoot[C]{\color{structurecolor}\small\thepage}%
  \renewcommand{\headrulewidth}{0pt}%
  \renewcommand{\footrulewidth}{0pt}%
}
